%Chargement des paquets
\usepackage{amsmath}
\usepackage{amsfonts}
\usepackage{amssymb}
\usepackage{enumerate}
\usepackage{mathtools}
\usepackage{bbm}
\usepackage{xparse, etoolbox}
\usepackage{enumerate}
\usepackage{enumitem}
\usepackage{mathabx}
\usepackage{babel}[french]

%environment
\usepackage{amsthm}
\usepackage{keytheorems}
\usepackage{hyperref} 

%Interface théorème
\newkeytheorem{lem}[
	name = Lemme
]

\newkeytheorem{prop}[
	name = Proposition
]

\newkeytheorem{thm}[
	name = Théorème
]

\newkeytheorem{coro}[
	name = Corollaire
]

\renewcommand*{\proofname}{Démonstration}

\theoremstyle{definition}
\newtheorem{defn}{Définition}

\theoremstyle{definition}
\newtheorem*{nota}{Notation}

\theoremstyle{definition}
\newtheorem*{limi}{Liminaire}

\theoremstyle{remark}
\newtheorem{rmq}{Remarque}

\theoremstyle{plain}
\newtheorem*{fait}{Fait}


%Ensembles de nombres
\newcommand{\N} {\mathbb{N}}
\newcommand{\Ne}{\N^\ast}
\newcommand{\Z} {\mathbb{Z}}
\newcommand{\Q} {\mathbb{Q}}
\newcommand{\R} {\mathbb{R}}
\newcommand{\Rb}{\overline{\mathbb{R}}}
\newcommand{\Rp}{\R_+}
\newcommand{\Rm}{\R_-}
\newcommand{\K} {\mathbb{K}}
\newcommand{\Ket} {\mathbb{K^{\ast}}}
\newcommand{\Cx}{\mathbb{C}}

%Ensembles, tribus
\newcommand{\A}{\mathbb{A}}
\renewcommand{\P}{\mathcal{P}}
\newcommand{\T}{\mathcal{T}}
\newcommand{\F}{\mathcal{F}}
\newcommand{\C} {\mathcal{C}}
\newcommand{\ouv}{\mathcal{O}}
\newcommand{\B}{\mathcal{B}}
\newcommand{\M}{\mathcal{M}}
\newcommand{\etag}{\mathcal{E}}
\newcommand{\etagOp}{\etag^{+}(\Omega)}
\newcommand{\etagO}{\etag(\Omega)}
%%Lp avant quotientage
\newcommand{\Lic}{\mathcal{L}^1}
\newcommand{\Lpc}{\mathcal{L}^p}
\newcommand{\Linfc}{\mathcal{L}^{\infty}}
\newcommand{\Lifull}{\Lic(\Omega, \B, m)}
\newcommand{\Lpfull}{\Lpc(\Omega, \B, m)}
\newcommand{\Linffull}{\Linfc(\Omega, \B, m)}
%%Lp après quotientage
\newcommand{\Li}{L^1}
\newcommand{\Ld}{L^2}
\newcommand{\Lp}{L^p}
\newcommand{\Linf}{L^{\infty}}
\newcommand{\Listd}{\Li(\Omega, m)} 
\newcommand{\Ldstd}{\Ld(\Omega, m)} 
\newcommand{\Lpstd}{\Lp(\Omega, m)} 

%Quelques opérateurs
\DeclareMathOperator{\codim}{codim}
\DeclareMathOperator{\rg}{rg}
\DeclareMathOperator{\tr}{tr}
\DeclareMathOperator{\vect}{Vect}
\DeclareMathOperator{\ima}{Im}
\DeclareMathOperator{\id}{Id}
\DeclareMathOperator{\sign}{sign}
\DeclareMathOperator{\ext}{ext}
\newcommand{\ind}{\mathbbm{1}}
\newcommand*{\spr}[2]{\left(\,#1\,\middle|\,#2\,\right)}
\newcommand{\interior}[1]{%
  {\kern0pt#1}^{\mathrm{o}}%
}
\DeclareMathOperator{\card}{card}
\DeclareMathOperator{\ppcm}{ppcm}

%Commandes de pure flemme
\newcommand{\veps}{\varepsilon}
\newcommand{\intab}{\int_{a}^{b}}
\newcommand{\intR}{\int_{-\infty}^{\infty}}
\newcommand{\intinf}{\int_{- \infty}^{+ \infty}}
\newcommand{\equi}{\Leftrightarrow}
\newcommand{\Kev}{$\K$-espace vectoriel }
\newcommand{\Kevs}{$\K$-espaces vectoriels }
\newcommand{\Rev}{$\R$-espace vectoriel }
\newcommand{\Revs}{$\R$-espaces vectoriels }
\newcommand{\Cev}{$\Cx$-espace vectoriel }
\newcommand{\Cevs}{$\Cx$-espaces vectoriels }
\newcommand{\evn}{(E, \norm{.})}
\newcommand{\interf}[2]{[#1,#2]}
\newcommand{\limb}{\lim\limits_}
\newcommand{\lebn}{\lambda_n}
\newcommand{\pp}{\text{~presque partout}}
\newcommand{\derivp}[2]{\frac{\partial #1}{\partial #2}}
\newcommand{\famiu}{(u_i)_{i \in I}}
\newcommand{\sev}{\text{~s.e.v.}}
\newcommand{\Mnp}{M_{n,p}(\K)}
\newcommand{\Mn}{M_{n}(\K)}
\newcommand{\tq}{\text{~t.q.~}}
\newcommand{\mq}{~montrer~que~}
\newcommand{\et}{\quad \text{et} \quad}
\newcommand{\ou}{\text{~ou~}}
\newcommand{\Kx}{\K[X]}


%Commandes de gars chelou
\renewcommand{\subset}{\subseteq}

%Commande restriction, ty duckduckgo
\def\restr#1#2{\mathchoice
              {\setbox1\hbox{${\displaystyle #1}_{\scriptstyle #2}$}
              \restraux{#1}{#2}}
              {\setbox1\hbox{${\textstyle #1}_{\scriptstyle #2}$}
              \restraux{#1}{#2}}
              {\setbox1\hbox{${\scriptstyle #1}_{\scriptscriptstyle #2}$}
              \restraux{#1}{#2}}
              {\setbox1\hbox{${\scriptscriptstyle #1}_{\scriptscriptstyle #2}$}
              \restraux{#1}{#2}}}
\def\restraux#1#2{{#1\,\smash{\vrule height .8\ht1 depth .85\dp1}}_{\,#2}} 

%Quelques normes
\DeclarePairedDelimiter\abs{\lvert}{\rvert}
\DeclarePairedDelimiter\labs{\bigl\lvert}{\bigr\rvert}
\DeclarePairedDelimiter\norm{\lVert}{\rVert}
\DeclarePairedDelimiterXPP\normi[1]{}\lVert\rVert{_1}{\ifblank{#1}{\:\cdot\:}{#1}}
\DeclarePairedDelimiterXPP\normd[1]{}\lVert\rVert{_2}{\ifblank{#1}{\:\cdot\:}{#1}}
\DeclarePairedDelimiterXPP\normp[1]{}\lVert\rVert{_p}{\ifblank{#1}{\:\cdot\:}{#1}}
\DeclarePairedDelimiterXPP\normq[1]{}\lVert\rVert{_q}{\ifblank{#1}{\:\cdot\:}{#1}}
\DeclarePairedDelimiterXPP\norminf[1]{}\lVert\rVert{_\infty}{\ifblank{#1}{\:\cdot\:}{#1}}

%Bracket en angle
\DeclarePairedDelimiter\angl{\langle}{\rangle}

%Set, parenthèses
\newcommand{\set}[1]{\{ #1 \}}
\newcommand{\bigset}[1]{\bigl\{ #1 \bigr\}}
\newcommand{\Bigset}[1]{\Bigl\{ #1 \Bigr\}}
\newcommand{\bigpar}[1]{\bigl( #1 \bigr)}
\newcommand{\Bigpar}[1]{\Bigl( #1 \Bigr)}

%Opitmisation des opérateurs de comparaison
\DeclarePairedDelimiter\tio{o (}{)}
\DeclarePairedDelimiter\gro{O (}{)}


\pagestyle{fancy}

\fancyhead[C]{Co-trigonalisation}
\fancyhead[R]{}

\begin{document}

\underline{Source :} Rombalide - Algèbre - page 678

\underline{Leçons :}
\begin{itemize}
\item 
\end{itemize}

\begin{nota}
$E$ est un \Kev de dimension finie $n \geq 1$ sur un corps $\K$.
\end{nota}

\begin{lem}
Si un endomorphisme $f$ de $E$ est trigonalisable et si $F$ est un sous espace vectoriel de $E$ stable par $f$, 
alors la restriction de $f$ à $F$ est trigonalisable.
\end{lem}

\begin{proof}
Par calcul direct, on constate que le polynôme caractéristique de $f$ est scindé, à fortiori, le polynôme caractéristique de $\restr{f}{F}$
l'est, d'où le résultat.
\end{proof}

\begin{thm}
Si $(f_i)_{i \in I}$ est une famille d’endomorphismes trigonalisables de $E$ qui
commutent deux à deux ($I$ ayant au moins deux éléments), il existe alors
une base commune de trigonalisation dans E pour la famille $(f_i)_{i \in I}$.
\end{thm}

\begin{proof}
Le résultat est clair pour une famille d'homothéties. Supposons donc que les $f_i$ ne sont pas tous des homothéties. \\

Montrons tout d'abord qu'il existe un vecteur propre commun à tous les $f_i$ en raisonnant par récurrence sur la dimension de $E$.
Pour $n=1$, il n'y a rien à démontrer (on revient au cas d'une famille d'homothéties). 
Supposons maintenant le résultat vrai pour toutes dimensions jusqu'à $n-1 \geq 1$ et supposons que $E$ est de dimension $n$.
Par hypothèses, il existe un indice $j \in I$ pour lequel $f_j$ n'est pas une homothétie. Etant trigonalisable, cet endomrophisme admet
néanmoins au moins une valeur propre et le sous-espace propre :
\[ E_{j, \lambda} = \set { x \in E ; f_j(x) = \lambda x } , \]
est de dimension au moins 1 et au plus $n-1$ ($f_j$ n'est pas une homothétie). 
De plus, par commutativité, $E_{j, \lambda}$ est stable pour chaque $f_i$ et, en conséquence, la restriction, à $E_{j, \lambda}$,
de chaque endomorphisme $f_i$ est trigonalisable. 
Ainsi, par hypothèse de récurrence, il existe un vecteur propre commun à tous les $\restr{f_i}{E_{j, \lambda}}$ et donc commun à tous
les $f_i$. \\

Montrons maintenant qu'il existe une base commune de trigonalisation. On procède à nouveau par récurrence sur la dimension de $E$.
Le résultat est évident pour $n=1$, supposons le acquis pour toutes les dimensions de 1 à $n-1 \geq 1$ et montrons le pour $E$
de dimension $n$.
On note $e_1$ un vecteur propre commun à tous les endomoprhismes $f_i$ et on complète ce vecteur en une base $B=(e_1, \dots, e_n)$.
La matrice de chaque endomorphisme $f_i$ dans $E$ est de la forme :
\[ A_i = \begin{pmatrix} 
\lambda_i & \mu_i \\
0 & B_i
\end{pmatrix} . \]
De plus, comme le polynôme caractéristique de chaque $f_i, P_{f_i} = (\lambda_i-X) P_{B_i}(X)$ est scindé sur $\K$ (chaque $f_i$ étant
trigonalisable), il en est de même de chaque polynôme $P_{B_i}(X)$ et donc la matrice $B_i$ est trigonalisable.
De plus on a :
\[ A_i A_j =  \begin{pmatrix} 
\lambda_i & \mu_i \\
0 & B_i
\end{pmatrix} 
 \begin{pmatrix} 
\lambda_j & \mu_j \\
0 & B_j
\end{pmatrix} 
=
 \begin{pmatrix} 
\lambda_i \lambda_j & \lambda_i \mu_j + mu_i B_j \\
0 & B_i B_j
\end{pmatrix} 
 \]
et de même :
\[ A_j A_i =  \begin{pmatrix} 
\lambda_j \lambda_i & \lambda_j \mu_i + mu_j B_i \\
0 & B_j B_i
\end{pmatrix} .
 \]
Comme les matrices $A_i$ commutent deux à deux, on en déduit que les matrices $B_i$ commutent aussi deux à deux.
En d'autres termes, les matrices $B_i$ définissent des endomorphises trigonalisables qui commutent deux à deux.
En leur appliquant l'hypothèse de récurrence, on trouve une base $(e'_2, \dots, e'_n)$ de l'espace vectoriel $F = \vect(e_2, \dots, e_n)$
dans lesquels les matrices de ces endomorphismes sont triangulaires. 
A fortiori, les matrices des endomrophismes $f_i$ dans la base $B'=(e_1, e'_2, \dots, e'_n)$ sont triangulaires.
\end{proof}

\end{document}